% ============================================================================
% PyQUDA — 完整技术文档
% Python wrapper for QUDA — GPU-accelerated Lattice QCD
% ============================================================================
\documentclass[12pt,a4paper]{ctexart}

% ---------- Packages ----------
\usepackage[hmargin=2.5cm,vmargin=2.5cm]{geometry}
\usepackage{hyperref}
\usepackage{graphicx}
\usepackage{listings}
\usepackage{xcolor}
\usepackage{enumitem}
\usepackage{booktabs}
\usepackage{tabularx}
\usepackage{caption}
\usepackage{fancyhdr}
\usepackage{titling}
\usepackage{tcolorbox}
\usepackage{amsmath,amssymb}
\usepackage{fontspec}
\setmonofont{DejaVu Sans Mono}[Scale=0.88]

% ---------- Code Listing Style ----------
\definecolor{codebg}{RGB}{245,245,245}
\definecolor{codeframe}{RGB}{200,200,200}
\definecolor{keyword}{RGB}{0,0,180}
\definecolor{string}{RGB}{0,128,0}
\definecolor{comment}{RGB}{128,128,128}

\lstdefinestyle{pycode}{
    backgroundcolor=\color{codebg}, frame=single, rulecolor=\color{codeframe},
    basicstyle=\footnotesize\ttfamily, keywordstyle=\color{keyword}\bfseries,
    stringstyle=\color{string}, commentstyle=\color{comment}\itshape,
    breaklines=true, showstringspaces=false, numbers=left,
    numberstyle=\tiny\color{gray}, numbersep=5pt, tabsize=4, language=Python,
}
\lstdefinestyle{bashcode}{
    backgroundcolor=\color{codebg}, frame=single, rulecolor=\color{codeframe},
    basicstyle=\footnotesize\ttfamily, breaklines=true, showstringspaces=false,
    numbers=left, numberstyle=\tiny\color{gray}, numbersep=5pt, language=bash,
}

\tcbuselibrary{skins,breakable}

% ---------- Metadata ----------
\title{{\Huge\bfseries PyQUDA}\\[4pt]
       {\large 完整技术文档}\\[4pt]
       {\normalsize Python Wrapper for QUDA — GPU-Accelerated Lattice QCD}}
\author{}
\date{2026年7月20日}

\pagestyle{fancy}
\fancyhf{}
\fancyhead[L]{\small PyQUDA 技术文档}
\fancyhead[R]{\small \thepage}
\renewcommand{\headrulewidth}{0.4pt}

\begin{document}

\maketitle
\thispagestyle{empty}

\vspace{2cm}
\begin{center}
\begin{tcolorbox}[width=0.85\textwidth,boxrule=0.8pt,colback=blue!3!white,colframe=blue!40!black]
\centering
\textbf{项目地址：}\url{https://github.com/CLQCD/PyQUDA.git}

\textbf{许可证：}MIT

\textbf{Python 版本要求：}3.8+

\textbf{当前版本：}0.10.53
\end{tcolorbox}
\end{center}

\newpage
\tableofcontents
\newpage

% ============================================================================
\section{项目概述}
% ============================================================================

PyQUDA 是 QUDA\footnote{QUDA: QCD on CUDA — 基于 CUDA/HIP/SYCL 的格点 QCD GPU 库} 的 Python 封装，使用 Cython 编写。它使格点 QCD 研究者能够从 Python 中利用 GPU 加速进行高性能计算，同时受益于 CuPy/PyTorch/NumPy 等优化数值库。

\subsection{设计目标}

\begin{itemize}[itemsep=3pt]
    \item \textbf{Python 优先}：所有 API 均为 Pythonic 风格，支持 NumPy/CuPy/PyTorch 数组互操作
    \item \textbf{多 GPU 支持}：通过 mpi4py 实现 MPI 分布式计算
    \item \textbf{多后端}：支持 CUDA、ROCm/HIP、SYCL（Intel GPU）三种 GPU 后端
    \item \textbf{完整功能}：覆盖夸克传播子、HMC、规范固定、规范 smearing、梯度流等
\end{itemize}

\subsection{核心功能}

\begin{table}[h]
\centering
\begin{tabular}{@{}ll@{}}
\toprule
\textbf{功能类别} & \textbf{具体内容} \\
\midrule
文件 I/O  & Chroma (QIO/LIME)、MILC、KYU、XQCD、NERSC、OpenQCD、ILDG、NPY、HDF5 \\
夸克传播子 & 各向同性/异性 Wilson、Clover（含多网格）、HISQ \\
HMC       & 规范作用量 + Clover/HISQ 费米子作用量 \\
规范固定  & Landau 规范（over-relaxation 方法）\\
规范 smearing & 3D/4D APE、Stout、HYP \\
费米子 smearing & Gaussian（Wuppertal）\\
梯度流    & Wilson/Symanzik 流（含费米子伴随流）\\
多网格    & 用于 Wilson/Clover 费米子的几何/代数多网格预条件子 \\
插件系统  & 自动从 C 头文件生成 Cython 封装 \\
\bottomrule
\end{tabular}
\caption{PyQUDA 核心功能清单}
\end{table}

\subsection{包结构}

PyQUDA 分为三个 PyPI 包：
\begin{itemize}[itemsep=2pt]
    \item \textbf{PyQUDA}（\texttt{pyquda\_core/}）：核心 Cython 绑定，直接封装 QUDA C API
    \item \textbf{PyQUDA-Utils}（\texttt{pyquda\_utils/}）：高层工具函数（反演、源构造、相位因子、I/O 辅助）
    \item \textbf{PyQUDA-IO}（\texttt{pyquda\_io/}）：低级文件 I/O，支持 9 种格点数据格式
\end{itemize}

% ============================================================================
\section{安装指南}
% ============================================================================

\subsection{前置条件}

\begin{table}[h]
\centering
\begin{tabular}{@{}ll@{}}
\toprule
\textbf{组件} & \textbf{要求} \\
\midrule
Python              & $\ge$ 3.8 \\
QUDA                & 从源码编译安装（develop 分支，commit $\ge$ PR\#1489）\\
MPI                 & mpi4py \\
GPU 后端（任选一）  & CuPy（CUDA/ROCm）或 PyTorch 或 DPNP（Intel）\\
编译工具链          & Cython, C/C++ 编译器, CUDA/HIP/SYCL 编译器 \\
\bottomrule
\end{tabular}
\caption{前置条件}
\end{table}

\subsection{从 PyPI 安装（推荐）}

\begin{lstlisting}[style=bashcode]
export QUDA_PATH=/path/to/quda/build/usqcd
python3 -m pip install pyquda pyquda-utils
\end{lstlisting}

\subsection{从源码安装}

\begin{lstlisting}[style=bashcode]
export QUDA_PATH=/path/to/quda/build/usqcd
git clone --recursive https://github.com/CLQCD/PyQUDA.git
cd PyQUDA/pyquda_core
python3 -m pip install .
cd ..
python3 -m pip install .
\end{lstlisting}

\subsection{开发模式安装（就地构建）}

\begin{lstlisting}[style=bashcode]
export QUDA_PATH=/path/to/quda/build/usqcd
git clone --recursive https://github.com/CLQCD/PyQUDA.git
cd PyQUDA/pyquda_core
python3 setup.py build_ext --inplace
cd ..
python3 setup.py egg_info
ln -s pyquda_core/pyquda_comm ./
ln -s pyquda_core/pyquda ./
\end{lstlisting}

\textbf{注意}：使用开发模式时，需要在运行脚本前将项目根目录添加到 \texttt{sys.path}。

% ============================================================================
\section{初始化配置}
% ============================================================================

\subsection{基本初始化}

PyQUDA 通过 \texttt{pyquda.init()} 函数完成所有初始化：

\begin{lstlisting}[style=pycode]
from pyquda import init
from pyquda_utils import core

# 基本初始化（4D 格点：4x4x4x8）
init([1, 1, 1, 1], [4, 4, 4, 8], backend="cupy")
\end{lstlisting}

\subsection{完整初始化参数}

\begin{lstlisting}[style=pycode]
def init(
    grid_size: Optional[List[int]] = None,      # MPI 网格大小
    latt_size: Optional[List[int]] = None,      # 全局格点大小 [Lx, Ly, Lz, Lt]
    grid_map: GridMapType = "default",          # MPI 网格映射策略
    backend: BackendType = "cupy",              # 数组后端
    backend_target: str = quda_define.target(), # GPU 目标
    init_quda: bool = True,                     # 是否初始化 QUDA
    use_malloc_quda: bool = False,              # 使用 QUDA 自定义分配器
    *,
    resource_path: str = "",                    # QUDA 资源路径
    rank_verbosity: List[int] = [0],            # MPI rank 详细度
    enable_mps: bool = False,                   # 启用 MPS
    enable_gdr: bool = False,                   # 启用 GPU Direct RDMA
    enable_p2p: int = 3,                        # P2P 通信策略
    enable_tuning: bool = True,                 # 启用自动调优
    enable_device_memory_pool: bool = True,     # 启用设备内存池
    enable_pinned_memory_pool: bool = True,     # 启用锁页内存池
    enable_managed_memory: bool = False,        # 启用托管内存
    deterministic_reduce: bool = False,         # 确定性归约
    device_reset: bool = False,                 # 设备重置
    max_multi_rhs: int = 0,                     # 最大多右端量
    ...
)
\end{lstlisting}

\textbf{网格映射类型}（\texttt{grid\_map}）：
\begin{itemize}[itemsep=2pt]
    \item \texttt{"default"} — 默认映射
    \item \texttt{"t\_fastest"} — t 方向最快
    \item \texttt{"x\_fastest"} — x 方向最快
    \item \texttt{"minimize"} — 最小化通信面
    \item \texttt{"shared"} — 共享内存
    \item \texttt{"dist\_graph"} — 分布式图
\end{itemize}

\textbf{数组后端}（\texttt{backend}）：\texttt{"numpy"}, \texttt{"cupy"}, \texttt{"torch"}, \texttt{"dpnp"}

\textbf{精度预设}：
\begin{itemize}[itemsep=2pt]
    \item \texttt{"none"} — 双精度所有层级
    \item \texttt{"invert"} — 混合精度反演
    \item \texttt{"multishift"} — 多 shift 反演精度
    \item \texttt{"multigrid"} — 多网格精度配置
\end{itemize}

\subsection{QUDA 环境变量}

\begin{lstlisting}[style=bashcode]
export QUDA_RESOURCE_PATH=/path/to/tuning/cache
export QUDA_ENABLE_TUNING=1
export QUDA_ENABLE_DEVICE_MEMORY_POOL=1
export QUDA_ENABLE_PINNED_MEMORY_POOL=1
\end{lstlisting}

% ============================================================================
\section{格点场类型}
% ============================================================================

PyQUDA 的核心数据结构是格点场（Lattice Field），存储在 GPU 显存中，具有特定的 even-odd 排序。

\subsection{LatticeInfo —— 格点几何信息}

\begin{lstlisting}[style=pycode]
class LatticeInfo:
    """
    描述格点几何的参数容器。

    Nd = 4         # 时空维度
    Ns = 4         # 自旋分量数
    Nc = 3         # 颜色数
    """
    def __init__(self, global_size: List[int],
                 t_boundary: int = 1,
                 anisotropy: float = 1.0): ...

    # 属性
    global_size: List[int]     # 全局格点大小 [Lx, Ly, Lz, Lt]
    size: List[int]            # 本地格点大小
    grid_size: List[int]       # MPI 网格大小
    grid_coord: List[int]      # 当前 rank 坐标
    volume: int                # 本地格点总体积 (Lx*Ly*Lz*Lt // 2)
    GLt: int                   # 全局时间方向大小
    Lt: int                    # 本地时间方向大小
\end{lstlisting}

\subsection{规范场}

\begin{lstlisting}[style=pycode]
class LatticeLink(latt_info, value=None):
    """单方向 SU(3) 链接变量
    形状: [2, Lt, Lz, Ly, Lx//2, Nc, Nc]"""

class LatticeGauge(latt_info, L5=4, value=None):
    """完整规范场（4 个方向）
    形状: [4, 2, Lt, Lz, Ly, Lx//2, Nc, Nc]"""

    # 核心方法
    def use()                     # 加载到 QUDA 并返回句柄
    def pack() / unpack()         # 打包/解包（用于通信）
    def covDev(x, mu)             # 协变导数
    def laplace(x, laplace3D)     # 拉普拉斯算子

    # Smearing
    def apeSmear(n_steps, alpha, ...)
    def stoutSmear(n_steps, rho, ...)
    def hypSmear(n_steps, alpha0, alpha1, alpha2, ...)

    # 梯度流
    def wilsonFlow(flow_time, n_steps, ...)
    def symanzikFlow(flow_time, n_steps, ...)
    def wilsonFlowScale(flow_time, n_steps, ...)

    # 可观测量
    def plaquette()               # 平均 plaquette
    def polyakovLoop()            # Polyakov 圈
    def energy()                  # 能量密度
    def qcharge() / qchargeDensity()  # 拓扑荷

    # 规范固定
    def fixingOVR(n_steps, tol, ...)   # Over-relaxation
    def fixingFFT(n_steps, tol, ...)   # FFT 加速

    # 随机数
    def gauss(seed)               # 高斯随机规范场

    # HDF5 I/O
    def loadH5(path) / saveH5(path) / appendH5(path) / updateH5(path)

class LatticeMom(latt_info, L5=4):
    """动量场（SU(3) 李代数）"""

class LatticeRotation(latt_info):
    """规范旋转场（用于规范固定）"""

class LatticeClover(latt_info):
    """Clover 项场"""
\end{lstlisting}

\subsection{费米子场（Wilson 型）}

\begin{lstlisting}[style=pycode]
class LatticeFermion(latt_info):
    """单个费米子场
    形状: [2, Lt, Lz, Ly, Lx//2, Ns, Nc]"""

class HalfLatticeFermion(latt_info):
    """Even/odd 半费米子场"""

class MultiLatticeFermion(latt_info, L5):
    """多右端量费米子场
    形状: [L5, 2, Lt, Lz, Ly, Lx//2, Ns, Nc]"""

class MultiHalfLatticeFermion(latt_info, L5):
    """多右端量半费米子场"""

class LatticePropagator(latt_info):
    """完全传播子
    形状: [2, Lt, Lz, Ly, Lx//2, Ns, Ns, Nc, Nc]"""
    # 支持 @ 运算符用于 gamma 矩阵乘法
\end{lstlisting}

\subsection{费米子场（Staggered 型）}

\begin{lstlisting}[style=pycode]
class LatticeStaggeredFermion(latt_info):
    """Staggered 费米子场，形状: [2, Lt, Lz, Ly, Lx//2, Nc]"""

class HalfLatticeStaggeredFermion(latt_info): ...
class MultiLatticeStaggeredFermion(latt_info, L5): ...
class MultiHalfLatticeStaggeredFermion(latt_info, L5): ...

class LatticeStaggeredPropagator(latt_info):
    """Staggered 传播子，形状: [2, Lt, Lz, Ly, Lx//2, Nc, Nc]"""
\end{lstlisting}

\subsection{Even-Odd 与 Lexicographic 转换}

\begin{lstlisting}[style=pycode]
def lexico(data: NDArray, axes: List[int], dtype=None) -> NDArray:
    """Even-odd → Lexicographic 排序转换"""

def evenodd(data: NDArray, axes: List[int], dtype=None) -> NDArray:
    """Lexicographic → Even-odd 排序转换"""
\end{lstlisting}

% ============================================================================
\section{Gamma 矩阵代数}
% ============================================================================

\begin{lstlisting}[style=pycode]
from pyquda import gamma

class Gamma(index: int, factor: int = 1):
    """单 gamma 矩阵 (0-15 索引)"""
    # 运算符: @, +, -, *, /, T, H, D (Dirac 共轭)

class Projector(factors: list):
    """16 个 gamma 矩阵的线性组合 (16 元素列表)"""
    # 支持相同的运算符

# 预定义常量
gamma.gamma_0, gamma.gamma_1, ..., gamma.gamma_4

# Gamma 基表示
class DeGrandRossiMatrix: ...   # DeGrand-Rossi (手征) 基
class DiracPauliMatrix: ...     # Dirac-Pauli 基
class DeGrandRossiSparse: ...   # 稀疏表示
class DiracPauliSparse: ...     # 稀疏表示
\end{lstlisting}

% ============================================================================
\section{Dirac 算子}
% ============================================================================

\subsection{类层次结构}

\begin{center}
\begin{tcolorbox}[width=0.85\textwidth,boxrule=0.5pt,colback=gray!5,colframe=gray!50]
\begin{verbatim}
Dirac (ABC)
├── GaugeDirac          —— 纯规范操作
└── FermionDirac (ABC)  —— 费米子 Dirac 算子
    ├── WilsonDirac
    ├── CloverWilsonDirac
    └── StaggeredFermionDirac (ABC)
        ├── StaggeredDirac
        └── HISQDirac
\end{verbatim}
\end{tcolorbox}
\end{center}

\subsection{GaugeDirac}

\begin{lstlisting}[style=pycode]
class GaugeDirac(latt_info):
    """纯规范操作算子"""
    def loadGauge(gauge)       # 加载规范场
    def freeGauge()            # 释放规范场
    def covDev(x, mu)          # 协变导数
    def laplace(x, laplace3D)  # 拉普拉斯算子
    def wuppertalSmear(x, steps, alpha)  # Wuppertal smearing

    # Smearing
    def apeSmear(n_steps, alpha, ...)
    def stoutSmear(n_steps, rho, ...)
    def hypSmear(n_steps, alpha0, alpha1, alpha2, ...)

    # 梯度流
    def wilsonFlow(flow_time, n_steps, ...)
    def symanzikFlow(flow_time, n_steps, ...)

    # 可观测量
    def plaquette() / polyakovLoop() / energy()
    def qcharge() / qchargeDensity()

    # 规范固定
    def fixingOVR(n_steps, tol, ...)
    def fixingFFT(n_steps, tol, ...)

    # 随机数
    def gaussGauge(seed) / gaussMom(seed)
\end{lstlisting}

\subsection{FermionDirac}

\begin{lstlisting}[style=pycode]
class FermionDirac(latt_info, mass, tol, maxiter, multigrid=None):
    """费米子 Dirac 算子基类"""

    # 基本操作
    def invert(b) -> x                          # 求解 D * x = b
    def mat(x) -> b                             # 计算 D * x
    def matDagMat(x) -> b                       # 计算 D^\dagger D * x
    def dslash(x, parity) -> b                  # 单 parity Dslash

    # 多右端量版本
    def invertMultiSrc(b) / matMultiSrc(x) / matDagMatMultiSrc(x)
    def dslashMultiSrc(x, parity)

    # 预条件版本
    def invertPC(b)                             # 预条件反演
    def invertCloverPC(b)                       # Clover 预条件反演

    # 配置
    def useGauge(gauge)                         # 上下文管理器加规范场
    def setPrecision(precision)                 # 设置精度
    def setReconstruct(reconstruct)             # 设置重构方案
    def setVerbosity(verbosity)                 # 设置详细度
\end{lstlisting}

\subsection{具体 Dirac 算子}

\begin{lstlisting}[style=pycode]
class WilsonDirac(latt_info, mass, tol, maxiter, multigrid=None):
    """Wilson 费米子 Dirac 算子"""

class CloverWilsonDirac(latt_info, mass, tol, maxiter,
                         clover_csw=0.0, clover_xi=1.0, multigrid=None):
    """Clover 改进 Wilson 费米子 Dirac 算子"""

class StaggeredDirac(latt_info, mass, tol, maxiter,
                      tadpole_coeff=1.0, multigrid=None):
    """Staggered 费米子 Dirac 算子"""

class HISQDirac(latt_info, mass, tol, maxiter,
                 naik_epsilon=0.0, multigrid=None):
    """HISQ 费米子 Dirac 算子"""
\end{lstlisting}

\subsection{多网格预条件子}

\begin{lstlisting}[style=pycode]
class Multigrid(param, inv_param, instance):
    """多网格预条件子生命周期管理"""
    def setParam(coarse_tol, coarse_maxiter,
                 setup_tol, setup_maxiter,
                 smoother_nu_pre, smoother_nu_post, smoother_omega)
    def new()                                  # 初始化
    def update(thin_update_only)               # 更新
    def destroy()                              # 销毁
\end{lstlisting}

% ============================================================================
\section{HMC（混合蒙特卡洛）}
% ============================================================================

\subsection{作用量类}

\begin{lstlisting}[style=pycode]
from pyquda.action import (
    GaugeAction, CloverWilsonAction, HISQAction,
    MultiHISQAction, LoopParam, RationalParam
)

# 规范作用量
gauge_action = GaugeAction(latt_info, loop_param, beta)

# Clover 费米子作用量（1 或 2 味）
clover_action = CloverWilsonAction(latt_info, rational_param,
    mass, n_flavor, tol, maxiter, clover_csw)

# HISQ 费米子作用量
hisq_action = HISQAction(latt_info, rational_param,
    tol, maxiter, naik_epsilon)

# 多味 HISQ（共享 fat/long 链接）
multi_hisq = MultiHISQAction(latt_info, [hisq1, hisq2, ...])
\end{lstlisting}

\textbf{LoopParam} 和 \textbf{RationalParam}：
\begin{lstlisting}[style=pycode]
LoopParam(path: List[List[int]], coeff: List[float])
    """规范圈参数 (如 plaquette [1,0,0,0], rectangle [2,0,0,0])"""

RationalParam(norm_force, residue_force, offset_force,
              norm_sample, residue_sample, offset_sample,
              norm_action, residue_action, offset_action)
    """有理近似参数（用于 RHMC）"""
\end{lstlisting}

\subsection{积分器类}

\begin{lstlisting}[style=pycode]
class Integrator(n_steps):
    """积分器基类"""

class INT_3G1F(n_steps):
    """3 阶规范 + 1 阶费米子"""

class O2Nf1Ng0V(n_steps):    # 2 阶, 1 道费米子力 (速度版本)
class O2Nf1Ng0P(n_steps):    # 2 阶, 1 道费米子力 (位置版本)
class O2Nf2Ng0V(n_steps):    # 2 阶, 2 道费米子力
class O2Nf2Ng0P(n_steps):    # 2 阶, 2 道费米子力
class O4Nf5Ng0V(n_steps):    # 4 阶, 5 道费米子力
class O4Nf5Ng0P(n_steps):    # 4 阶, 5 道费米子力
\end{lstlisting}

\subsection{HMC 类}

\begin{lstlisting}[style=pycode]
class HMC(latt_info, monomials, integrator, hmc_inner=None):
    """混合蒙特卡洛"""

    def initialize(seed, gauge, mom=None):
        """初始化轨迹"""

    # 动力学
    def gaussMom()                    # 生成高斯随机动量
    def integrate(t, project_tol)     # 执行轨迹积分
    def accept(delta_s) -> bool       # Metropolis 接受/拒绝

    # 作用量
    def momAction()        # 动量作用量
    def gaugeAction()      # 规范作用量
    def fermionAction(use_force_param)  # 费米子作用量

    # I/O
    def loadGauge() / saveGauge(path)
    def loadMom() / saveMom(path)

    # 可观测量
    def plaquette() / polyakovLoop() / energy()

    # 伪费米子
    def samplePhi()        # 生成伪费米子

    # HDF5
    def loadH5(path) / saveH5(path)
\end{lstlisting}

% ============================================================================
\section{工具函数（pyquda\_utils）}
% ============================================================================

\subsection{核心工具（core.py）}

\begin{lstlisting}[style=pycode]
# 传播子反演
def invert(dirac, source_type, t_srce, source_phase, mrhs, restart) -> LatticePropagator
def invertEigenvector(dirac, t_srce, source_propag, mrhs, restart)
def invertSequential(dirac, source_propag, t_srce, mrhs, restart)
def invertPropagator(dirac, source_propag, mrhs, restart)

# Staggered 版本
def invertStaggered(dirac, source_type, t_srce, source_phase, mrhs, restart)
def invertStaggeredSequential(...)
def invertStaggeredPropagator(...)

# MPI 数据收集/分发
def gatherLattice(data, axes, reduce_op, root)
def gatherLattice2(data, tzyx, reduce_op, root)
def scatterLattice(data_all, tzyx, root)
def gatherScatterLattice(data, tzyx, reduce_op, root)

# Dirac 算子工厂
def getDirac(latt_info, mass, tol, maxiter, xi_0, clover_coeff_t,
             clover_coeff_r, multigrid) -> FermionDirac
def getClover(...) -> CloverWilsonDirac
def getWilson(...) -> WilsonDirac
def getStaggered(...) -> StaggeredDirac
def getHISQ(...) -> HISQDirac
\end{lstlisting}

\subsection{源构造（source.py）}

\begin{lstlisting}[style=pycode]
def point(latt_info, t_srce, spin, color, phase) -> point 源
def wall(latt_info, t_srce, spin, color, phase) -> wall 源
def volume(latt_info, spin, color, phase) -> volume 源
def momentum(latt_info, t_srce, spin, color, phase) -> 动量源
def colorvector(latt_info, t_srce, spin, phase) -> 颜色向量源
def source(latt_info, source_type, t_srce, spin, color, source_phase) -> 调度器

# 便捷函数（自动生成全自旋/全颜色组合）
def source12(...)  # 12 个 spin-color 源 → LatticePropagator
def source3(...)   # 3 个 color 源 → LatticeStaggeredPropagator

# Gaussian smearing
def gaussianSmear(x, gauge, rho, n_steps)
def gaussian(x, gauge, rho, n_steps) / gaussian12(...) / gaussian3(...)

# Sequential 传播子
def sequential(x, t_srce) / sequential12(...) / sequential3(...)
\end{lstlisting}

\subsection{相位因子（phase.py / phase\_v2.py）}

\begin{lstlisting}[style=pycode]
# 动量列表生成
def getMomList(mom2_max, mom2_min=0) -> list[tuple]
def getMomDict(mom2_max, mom2_min=0) -> dict[int, str]

# 相位计算
class MomentumPhase(latt_info):
    """GPU 上的动量相位因子"""

# v2 API
class LocationPhase(latt_info):  # 格点坐标相位
class DistancePhase(latt_info):  # 距离相位
class MomentumPhase(latt_info):
    def getPhase(mom, x0)           # 单动量相位
    def getPhases(mom_mode_list, x0) # 多动量批量
\end{lstlisting}

\subsection{Gamma 工具（gamma.py）}

\begin{lstlisting}[style=pycode]
def gamma_mul_fermion(gamma, fermion)           # γ × fermion
def propagator_mul_gamma(propagator, gamma)     # propagator × γ
def gamma_mul_propagator(gamma, propagator)     # γ × propagator
def gamma_mul_propagator_mul_gamma(g1, prop, g2) # γ₁ × propagator × γ₂

# Monkey-patch 的运算符支持
LatticeFermion @ Gamma       # gamma_mul_fermion
LatticePropagator @ Gamma    # propagator_mul_gamma
Gamma @ LatticePropagator    # gamma_mul_propagator
\end{lstlisting}

\subsection{其他工具}

\begin{lstlisting}[style=pycode]
# 校验和
def checksumSciDAC(field) -> (sum29, sum31)

# 场类型转换
def linkToFermion(link) / fermionToLink(fermion)
def multiFermionToPropagator(multi_fermion)
def propagatorToMultiFermion(propagator)

# Wilson 圈
def wilson_loop(gauge, path) -> LatticeLink

# 3D Laplace 本征求解
def laplace3d(gauge, n_ev, n_kr, tol, max_restarts, poly_deg, poly_cut) -> (evals, evecs)

# FFT（MPI 分布式）
def transform(latt_info, field_shape, recvdim, senddim, buf)

# HMC 参数
def wilsonGaugeLoopParam() -> LoopParam
def symanzikTreeGaugeLoopParam(u_0) -> LoopParam
def symanzikOneLoopGaugeLoopParam(u_0, fermion_type, n_flavor) -> LoopParam
def fermionRationalParam(n_flavor, ...) -> RationalParam
def staggeredFermionRationalParam(masses, n_flavors, ...) -> RationalParam

# 多网格验证
def real_geo_block_size(sublatt_size, geo_block_size)
\end{lstlisting}

% ============================================================================
\section{文件 I/O}
% ============================================================================

PyQUDA 支持 9 种格点数据格式，分为低级 API（\texttt{pyquda\_io}）和高级 API（\texttt{pyquda\_utils.io}）。

\subsection{低级 I/O API}

\begin{table}[h]
\centering
\begin{tabular}{@{}llll@{}}
\toprule
\textbf{格式} & \textbf{模块} & \textbf{规范场} & \textbf{传播子} \\
\midrule
Chroma QIO (LIME) & \texttt{chroma.py} & 读 & 读（Wilson + Staggered）\\
MILC              & \texttt{milc.py}   & 读/写 & 读（Wilson + Staggered）\\
ILDG (LIME)       & \texttt{ildg.py}   & 读 & — \\
KYU               & \texttt{kyu.py}    & 读/写 & 读/写 \\
XQCD              & \texttt{xqcd.py}   & — & 读/写 \\
NERSC             & \texttt{nersc.py}  & 读/写 & — \\
OpenQCD           & \texttt{openqcd.py}& 读/写 & — \\
NPY (NumPy)       & \texttt{npy.py}    & 读/写 & 读/写 \\
\bottomrule
\end{tabular}
\caption{低级 I/O 格式支持}
\end{table}

\subsection{高级 I/O API（返回 Field 对象）}

\begin{lstlisting}[style=pycode]
from pyquda_utils.io import (
    # Chroma QIO
    readChromaQIOGauge, readChromaQIOPropagator, readChromaQIOStaggeredPropagator,

    # ILDG
    readILDGGauge, readILDGBinGauge,

    # MILC
    readMILCGauge, writeMILCGauge,
    readMILCQIOPropagator, readMILCQIOStaggeredPropagator,

    # KYU
    readKYUGauge, writeKYUGauge,
    readKYUPropagator, writeKYUPropagator,

    # XQCD
    readXQCDPropagator, writeXQCDPropagator,
    readXQCDStaggeredPropagator, writeXQCDStaggeredPropagator,

    # NPY
    readNPYGauge, writeNPYGauge,
    readNPYPropagator, writeNPYPropagator,

    # NERSC
    readNERSCGauge, writeNERSCGauge,

    # OpenQCD
    readOpenQCDGauge, writeOpenQCDGauge,

    # Gamma 基转换
    rotateToDiracPauli, rotateToDeGrandRossi,
)
\end{lstlisting}

\textbf{使用示例}：
\begin{lstlisting}[style=pycode]
from pyquda_utils.io import readChromaQIOGauge

# 读取 Chroma 格式规范场
gauge = readChromaQIOGauge(latt_info, "config.lime", "config.ini.xml")

# 读取 MILC 格式规范场
gauge = readMILCGauge(latt_info, "config.milc")
\end{lstlisting}

% ============================================================================
\section{插件系统}
% ============================================================================

PyQUDA 的插件系统使用 \texttt{pycparser} 自动从 C 头文件生成 Cython 封装。

\subsection{使用插件构建工具}

\begin{lstlisting}[style=bashcode]
# 自动生成封装
pyquda_plugins -l <库名> -i <头文件.h> -L <库路径> -I <包含路径>
\end{lstlisting}

\subsection{内置插件}

\subsubsection{pycontract —— 强子收缩}

\begin{lstlisting}[style=pycode]
from pyquda_plugins.pycontract import (
    # 介子
    mesonTwoPoint, mesonAllSinkTwoPoint, mesonAllSourceTwoPoint,

    # 重子
    baryonDiquark, baryonTwoPoint, baryonGeneralTwoPoint,
    baryonSequentialTwoPoint, baryonTwoPoint_v2,
)
\end{lstlisting}

\subsubsection{pygwu —— GWU 重叠费米子}

\begin{lstlisting}[style=pycode]
from pyquda_plugins.pygwu import Overlap

ov = Overlap(latt_info)
ov.buildHWilson(gauge, kappa)       # 构造 H_Wilson
ov.buildHWilsonEigen(eignum, ...)    # 计算本征向量
ov.buildOverlap(prec, use_fp32)     # 构造重叠算子
ov.invert(b, masses, tol, maxiter, ...)  # 多质量重叠反演
\end{lstlisting}

% ============================================================================
\section{示例演示}
% ============================================================================

\subsection{示例 1：淬火 HMC}

\begin{lstlisting}[style=pycode]
# 1_Quenched_HMC.py
from pyquda import init, LatticeInfo, LatticeGauge, LatticeMom
from pyquda.action import GaugeAction, LoopParam
from pyquda.hmc import HMC, O2Nf1Ng0V

# 初始化
latt_info = LatticeInfo([16, 16, 16, 32])
init([1, 1, 1, 1], [16, 16, 16, 32], backend="cupy")

# 定义作用量 (Wilson 规范作用量)
loop_param = LoopParam(path=[[1, 0, 0, 0]], coeff=[4.0 / 3])
gauge_action = GaugeAction(latt_info, loop_param, beta=5.8)

# 定义积分器和 HMC
integrator = O2Nf1Ng0V(n_steps=10)
hmc = HMC(latt_info, [gauge_action], integrator)

# 初始化并运行轨迹
gauge = LatticeGauge(latt_info)
gauge.gauss(seed=42)
hmc.initialize(seed=42, gauge=gauge)

for traj in range(100):
    hmc.gaussMom()
    delta_s = hmc.integrate(1.0, 1e-7)
    if hmc.accept(delta_s):
        print(f"Traj {traj}: accepted, dS={delta_s:.4f}")
        hmc.saveGauge(f"config_{traj}.npy")
\end{lstlisting}

\subsection{示例 2：夸克传播子}

\begin{lstlisting}[style=pycode]
# 2_Quark_Propagator.py
from pyquda import LatticeInfo, LatticeGauge, LatticeFermion
from pyquda.dirac import CloverWilsonDirac
from pyquda_utils.io import readChromaQIOGauge
from pyquda_utils.core import invert
from pyquda_utils.source import wall

# 初始化格点
latt_info = LatticeInfo([24, 24, 24, 72])
init([1, 1, 1, 1], [24, 24, 24, 72])

# 读取规范场并 smearing
gauge = readChromaQIOGauge(latt_info, "config.lime", "config.ini.xml")
gauge.stoutSmear(n_steps=1, rho=0.125)

# 创建 Dirac 算子并反演
dirac = CloverWilsonDirac(latt_info, mass=-0.2395, tol=1e-12, maxiter=10000)
dirac.setPrecision("invert")

with dirac.useGauge(gauge):
    prop = invert(dirac, wall, t_srce=0, source_phase=[0, 0, 0, 0], mrhs=12, restart=0)
\end{lstlisting}

\subsection{示例 3：Pion 和质子两点函数}

\begin{lstlisting}[style=pycode]
# 3_Pion_Proton_2pt.py
from pyquda_utils.source import source12
from pyquda_plugins.pycontract import mesonTwoPoint, baryonTwoPoint

# 12 个 spin-color 源 → 传播子
prop = source12(latt_info, "wall", t_srce=0, source_phase=[0, 0, 0, 0])

with dirac.useGauge(gauge):
    prop = invert(dirac, prop, restart=0)

# Pion 2pt: gamma5 投影
pion_2pt = mesonTwoPoint(prop, prop, gamma_ab=gamma_4,
                          gamma_dc=gamma_4)  # γ₄ × propagator × γ₄

# 质子 2pt: Cγ₅ 投影
proton_2pt = baryonTwoPoint(prop, prop, prop,
    contract_type="Cg5", gamma_mn=gamma_5)
\end{lstlisting}

\subsection{示例 4：Pion PCAC 关系}

\begin{lstlisting}[style=pycode]
# 4_Pion_PCAC.py
# 计算轴矢 Ward 恒等式
# ⟨∂_μ A_μ(x) O(0)⟩ = 2 m_q ⟨P(x) O(0)⟩

# A₄ 插入：γ₅γ₄
pion_a4 = mesonTwoPoint(prop_axial, prop, gamma_ab=gamma_4 * gamma_5,
                         gamma_dc=gamma_4 * gamma_5)

# 赝标量插入：γ₅
pion_p = mesonTwoPoint(prop, prop, gamma_ab=gamma_4 * gamma_5,
                        gamma_dc=gamma_4 * gamma_5)

# PCAC 质量: m_q = ⟨∂A₄⟩ / (2⟨P⟩)
\end{lstlisting}

\subsection{示例 5：Pion 色散关系}

\begin{lstlisting}[style=pycode]
# 5_Pion_Dispersion.py
from pyquda_utils.phase import MomentumPhase

phase = MomentumPhase(latt_info)

for mom_str in ["PX0PY0PZ0", "PX1PY0PZ0", "PX1PY1PZ0", "PX1PY1PZ1"]:
    mom_phase = phase.getPhase(mom_str, x0=0)
    # 动量投影的有效质量
    # meff(t) = arccosh((C(t+1) + C(t-1)) / (2 * C(t)))
\end{lstlisting}

% ============================================================================
\section{接口清单（API Reference）}
% ============================================================================

\subsection{pyquda 核心包}

\begin{lstlisting}[style=pycode]
# __init__.py —— 初始化
init(grid_size, latt_size, grid_map, backend, backend_target, init_quda, ...)
initQUDA(grid_size, device, use_malloc_quda)
isQUDAInitialized() -> bool

# field.py —— 格点场类型
LatticeInfo(global_size, t_boundary, anisotropy)
LatticeInt, MultiLatticeInt
LatticeReal, MultiLatticeReal
LatticeComplex, MultiLatticeComplex
LatticeLink, LatticeGauge, LatticeMom, LatticeRotation, LatticeClover
LatticeFermion, HalfLatticeFermion, MultiLatticeFermion, MultiHalfLatticeFermion
LatticePropagator
LatticeStaggeredFermion, HalfLatticeStaggeredFermion
MultiLatticeStaggeredFermion, MultiHalfLatticeStaggeredFermion
LatticeStaggeredPropagator
lexico(data, axes, dtype), evenodd(data, axes, dtype)

# gamma.py —— Gamma 代数
Gamma(index, factor)
Projector(factors)
DeGrandRossiMatrix, DiracPauliMatrix
DeGrandRossiSparse, DiracPauliSparse

# hmc.py —— HMC
Integrator(n_steps)
INT_3G1F, O2Nf1Ng0V, O2Nf1Ng0P, O2Nf2Ng0V, O2Nf2Ng0P
O4Nf5Ng0V, O4Nf5Ng0P
HMC(latt_info, monomials, integrator, hmc_inner)
\end{lstlisting}

\subsection{pyquda.action}

\begin{lstlisting}[style=pycode]
LoopParam(path, coeff)
RationalParam(norm_force, residue_force, offset_force, ...)
Action(latt_info, dirac)
GaugeAction(latt_info, loop_param, beta)
FermionAction(latt_info, dirac, rational_param, mass, n_flavor, tol, maxiter)
CloverWilsonAction(latt_info, rational_param, mass, n_flavor, tol, maxiter, clover_csw)
HISQAction(latt_info, rational_param, tol, maxiter, naik_epsilon)
MultiHISQAction(latt_info, hisq_actions)
\end{lstlisting}

\subsection{pyquda.dirac}

\begin{lstlisting}[style=pycode]
Dirac(latt_info)
GaugeDirac(latt_info)
    loadGauge, freeGauge, covDev, laplace, wuppertalSmear
    apeSmear, stoutSmear, hypSmear
    wilsonFlow, symanzikFlow, wilsonFlowScale
    plaquette, polyakovLoop, energy, qcharge, qchargeDensity
    fixingOVR, fixingFFT, projectSU3, staggeredPhase
    gaussGauge, gaussMom

FermionDirac(latt_info, mass, tol, maxiter, multigrid)
    invert, invertRestart, mat, matDagMat, dslash
    invertMultiSrc, invertMultiSrcRestart
    matMultiSrc, matDagMatMultiSrc, dslashMultiSrc
    invertPC, invertCloverPC, useGauge
    setPrecision, setReconstruct, setVerbosity

WilsonDirac(latt_info, mass, tol, maxiter, multigrid)
CloverWilsonDirac(latt_info, mass, tol, maxiter, clover_csw, clover_xi, multigrid)
StaggeredDirac(latt_info, mass, tol, maxiter, tadpole_coeff, multigrid)
HISQDirac(latt_info, mass, tol, maxiter, naik_epsilon, multigrid)

Multigrid(param, inv_param, instance)
    setParam, new, update, destroy
Deflation(param, inv_param, instance)
    new, update, destroy
\end{lstlisting}

\subsection{pyquda\_utils.core}

\begin{lstlisting}[style=pycode]
init(...)
invert(dirac, source_type, t_srce, source_phase, mrhs, restart)
invertEigenvector(dirac, t_srce, source_propag, mrhs, restart)
invertSequential(dirac, source_propag, t_srce, mrhs, restart)
invertPropagator(dirac, source_propag, mrhs, restart)
invertStaggered(dirac, source_type, t_srce, source_phase, mrhs, restart)
invertStaggeredSequential(...)
invertStaggeredPropagator(...)
gatherLattice(data, axes, reduce_op, root)
gatherLattice2(data, tzyx, reduce_op, root)
scatterLattice(data_all, tzyx, root)
gatherScatterLattice(data, tzyx, reduce_op, root)
getDirac(latt_info, mass, tol, maxiter, xi_0, clover_coeff_t, clover_coeff_r, multigrid)
getClover(...)
getWilson(...)
getStaggered(...)
getHISQ(...)
\end{lstlisting}

\subsection{pyquda\_utils.source}

\begin{lstlisting}[style=pycode]
point(latt_info, t_srce, spin, color, phase)
wall(latt_info, t_srce, spin, color, phase)
volume(latt_info, spin, color, phase)
momentum(latt_info, t_srce, spin, color, phase)
colorvector(latt_info, t_srce, spin, phase)
source(latt_info, source_type, t_srce, spin, color, source_phase)
source12(latt_info, source_type, t_srce, source_phase)
source3(latt_info, source_type, t_srce, source_phase)
gaussianSmear(x, gauge, rho, n_steps)
gaussian(x, gauge, rho, n_steps)
gaussian12(x12, gauge, rho, n_steps)
gaussian3(x3, gauge, rho, n_steps)
sequential(x, t_srce)
sequential12(x12, t_srce)
sequential3(x3, t_srce)
\end{lstlisting}

\subsection{pyquda\_utils.phase / phase\_v2}

\begin{lstlisting}[style=pycode]
# phase.py
getMomList(mom2_max, mom2_min=0)
getMomDict(mom2_max, mom2_min=0)
MomentumPhase(latt_info)

# phase_v2.py
LocationPhase(latt_info)
DistancePhase(latt_info)
MomentumPhase(latt_info): getPhase(mom, x0), getPhases(mom_list, x0)
\end{lstlisting}

\subsection{pyquda\_utils 其他工具}

\begin{lstlisting}[style=pycode]
# gamma.py
gamma_mul_fermion(gamma, fermion)
propagator_mul_gamma(propagator, gamma)
gamma_mul_propagator(gamma, propagator)
gamma_mul_propagator_mul_gamma(g1, prop, g2)

# checksum.py
checksumSciDAC(field) -> (sum29, sum31)

# convert.py
linkToFermion(link), fermionToLink(fermion)
multiFermionToPropagator(multi_fermion)
propagatorToMultiFermion(propagator)
multiField(fields)

# wilson_loop.py
wilson_loop(gauge, path) -> LatticeLink

# eigensolve.py
laplace3d(gauge, n_ev, n_kr, tol, max_restarts, poly_deg, poly_cut)

# fft.py
transform(latt_info, field_shape, recvdim, senddim, buf)
loadBalanceSubsize(GL, size, rank), loadBalanceStart(GL, size, rank)

# hmc_param.py
wilsonGaugeLoopParam(), symanzikTreeGaugeLoopParam(u_0)
symanzikOneLoopGaugeLoopParam(u_0, fermion_type, n_flavor)
fermionRationalParam(n_flavor, ...)
staggeredFermionRationalParam(masses, n_flavors, ...)

# geo_block_size.py
real_geo_block_size(sublatt_size, geo_block_size)

# deprecated.py
cb2(), smear(), smear4(), invert12(), getDslash(), getStaggeredDslash()
\end{lstlisting}

\subsection{pyquda\_io —— 低级 I/O}

\begin{lstlisting}[style=pycode]
# io_general.py
IOGeneral, readIOGeneral, writeIOGeneral

# chroma.py
readQIOGauge, readQIOPropagator, readQIOStaggeredPropagator

# milc.py
readGauge, writeGauge, readQIOPropagator, readQIOStaggeredPropagator

# ildg.py
readGauge, readBinGauge

# kyu.py
readGauge, writeGauge, readPropagator, writePropagator

# xqcd.py
readPropagator, writePropagator, readStaggeredPropagator, writeStaggeredPropagator
readPropagatorFast, writePropagatorFast

# nersc.py
readGauge, writeGauge

# openqcd.py
readGauge, writeGauge

# npy.py
readGauge, writeGauge, readPropagator, writePropagator
\end{lstlisting}

\subsection{pyquda\_plugins}

\begin{lstlisting}[style=pycode]
# pycontract
init(), mesonTwoPoint, mesonAllSinkTwoPoint, mesonAllSourceTwoPoint
baryonDiquark, baryonTwoPoint, baryonGeneralTwoPoint
baryonSequentialTwoPoint, baryonTwoPoint_v2

# pygwu
Overlap(latt_info):
    buildHWilson, loadHWilsonEigen, buildHWilsonEigen
    buildOverlap, loadOverlapEigen
    invert(b, masses, tol, maxiter, ...)
readEigenSystem, readEigenSystemInChunks
multiFermionFromDiracPauli, multiFermionToDiracPauli, ...
\end{lstlisting}

% ============================================================================
\section{完整示例：端到端介子谱计算}
% ============================================================================

以下示例整合了 PyQUDA 的核心流程，从读取规范场到提取有效质量：

\begin{lstlisting}[style=pycode]
from pyquda import init, LatticeInfo, LatticeGauge
from pyquda.dirac import CloverWilsonDirac
from pyquda_utils.core import invert
from pyquda_utils.source import source12
from pyquda_utils.io import readChromaQIOGauge
from pyquda_plugins.pycontract import mesonTwoPoint
from pyquda_utils.phase import MomentumPhase
import numpy as np

# 1. 初始化
latt_info = LatticeInfo([24, 24, 24, 72])
init([1, 1, 1, 1], [24, 24, 24, 72], backend="cupy")

# 2. 读取并处理规范场
gauge = readChromaQIOGauge(latt_info, "config.lime", "config.ini.xml")
gauge.stoutSmear(n_steps=1, rho=0.125)

# 3. 创建 Dirac 算子
dirac = CloverWilsonDirac(latt_info, mass=-0.2395, tol=1e-12,
                           maxiter=10000, clover_csw=1.0)
dirac.setPrecision("invert")

# 4. 计算传播子
prop = source12(latt_info, "wall", t_srce=0,
                 source_phase=[0, 0, 0, 0])

with dirac.useGauge(gauge):
    prop = invert(dirac, prop, restart=0)

# 5. 计算 pion 两点函数
from pyquda import gamma
pion_2pt = mesonTwoPoint(prop, prop,
    gamma_ab=gamma.gamma_4 * gamma.gamma_5,
    gamma_dc=gamma.gamma_4 * gamma.gamma_5)

# 6. 提取有效质量
data = pion_2pt.data  # numpy/cupy 数组
meff = np.arccosh((data[:-2] + data[2:]) / (2 * data[1:-1]))
print("Effective mass:", meff)
\end{lstlisting}

% ============================================================================
\section{依赖关系}
% ============================================================================

\begin{table}[h]
\centering
\begin{tabular}{@{}lll@{}}
\toprule
\textbf{依赖} & \textbf{版本} & \textbf{用途} \\
\midrule
QUDA          & develop 分支  & GPU 加速格点 QCD 核心库 \\
Python        & $\ge$ 3.8     & 运行环境 \\
mpi4py        & —             & MPI 分布式计算 \\
numpy         & —             & 数值计算（全后端兼容层）\\
Cython        & —             & C/CUDA 绑定生成 \\
\multicolumn{3}{@{}l}{\textbf{可选 GPU 后端（至少一个）}} \\
CuPy          & —             & CUDA/ROCm GPU 数组 \\
PyTorch       & —             & CUDA 张量计算 \\
DPNP          & —             & Intel GPU 数组 \\
\multicolumn{3}{@{}l}{\textbf{可选工具}} \\
h5py          & —             & HDF5 I/O \\
opt\_einsum   & —             & 张量收缩优化 \\
gmpy2         & —             & 有理近似 (Remez 算法) \\
matplotlib    & —             & 绑图（示例）\\
GPT           & —             & Grid Python Toolkit 互操作 \\
\bottomrule
\end{tabular}
\caption{完整依赖列表}
\end{table}

% ============================================================================
\section{维护与贡献}
% ============================================================================

\subsection{手动更新注意事项}

以下文件需要手动更新（无法自动生成）：
\begin{itemize}[itemsep=2pt]
    \item \texttt{quda.in.pxd} 和 \texttt{quda.in.pyx} 中的函数定义
    \item \texttt{pyquda.pyi} 和 \texttt{enum\_quda.in.py} 中的 docstring
\end{itemize}

这意味着只要 QUDA API 保持不变，PyQUDA 就能与未来版本兼容。

\subsection{开发环境建议}

\begin{itemize}[itemsep=2pt]
    \item 使用就地构建模式（\texttt{setup.py build\_ext --inplace}）而非安装模式
    \item 修改 Python 文件后立即生效，无需重新编译
    \item 在集群上运行时需要将项目根目录添加到 \texttt{sys.path}
\end{itemize}

\end{document}
